\section{Operaciones en Cuaterniones}
Podemos aplicar operaciones en los cuaterniones. Podemos sumar, multiplicar (equivalente al producto cruz) y aplicar el producto punto entre cuaterniones. También podemos calcular la norma y obtener el conjuntado de un cuaternión. 

No se cumple la conmutatividad en el producto de cuaterniones. Sin embargo, sí se cumple la distributividad. 

En las siguientes subsecciones, veremos cómo aplicar estas operaciones.

\subsection{Operaciones entre las Componentes Imaginarias}
Para poder aplicar operaciones en cuaterniones, primero debemos saber cómo operar entre las componentes imaginarias. Es decir, necesitamos saber qué pasa cuando multipliquemos por ejemplo $ij$ o $ji$ o $kj$. 

Las respuestas se obtienen de la relación:

\begin{equation}
    i^{2}=j^{2}=k^{2}=ijk=-1
\end{equation}

e.g. Sean 

\subsection{Suma}

\subsection{Producto}

\subsection{Producto punto}

\subsection{Norma}

\subsection{Conjugado}

\subsection{No Conmutatividad}

\subsection{Distributividad}





